\documentclass{article}

\usepackage[final]{neurips_2019}
\usepackage[utf8]{inputenc} 
\usepackage[T1]{fontenc}    
\usepackage{hyperref}       
\usepackage{url}            
\usepackage{booktabs}       
\usepackage{amsfonts}       
\usepackage{nicefrac}       
\usepackage{microtype}      
\usepackage{graphicx}
\usepackage{amsmath}

\title{Midterm Report: [Your Project Title Here]}

\author{%
  Tianxiang Wang \\
  Information Networking Institute\\
  Carnegie Mellon University\\
  Pittsburgh, PA 15213 \\
  \texttt{jerry945666@gmail.com} \\
  \And
  [Partner Name] \\
  [Department] \\
  Carnegie Mellon University \\
  \texttt{[email]@andrew.cmu.edu} \\
}

\begin{document}

\maketitle

\begin{abstract}
  % Rubric: The problem statement is understandable to non-experts.
  The problem statement is understandable to non-experts.
\end{abstract}

\section{Introduction}
\subsection{Motivation and Objectives}
% Rubric: Students articulate why the problem is important or who benefits 
% and state at least one qualitative or quantitative objective.
Students articulate why the problem is important or who benefits and state at least one qualitative or quantitative objective.

\section{Related Work and Background}
\subsection{Literature Review}
% Rubric: Students cite 1–2 key papers related to their project, 
% with brief connections to the problem.
Students cite 1–2 key papers related to their project, with brief connections to the problem.

\subsection{Background}
% Rubric: At least one cited paper connects the evaluation metric or baseline to the problem.
At least one cited paper connects the evaluation metric or baseline to the problem.

\section{Methodology}
\subsection{Model Description}
% Rubric: Basic explanation of the proposed model (inputs/outputs). 
% Diagrams/tables encouraged.
Basic explanation of the proposed model (inputs/outputs). Diagrams/tables encouraged.

\subsection{Dataset}
% Rubric: Data source is identified and basic preparation steps outlined.
Data source is identified and basic preparation steps outlined.

\subsection{Evaluation Metric}
% Rubric: The chosen metric is stated, with a simple explanation of its relevance.
The chosen metric is stated, with a simple explanation of its relevance.

\subsection{Loss Function}
% Rubric: Mentioned briefly if applicable.
Mentioned briefly if applicable.

\subsection{Experimental Depth and Iteration}
% Rubric: Students present evidence of initial iterations 
% (e.g., at least two runs, adjustments, or configurations attempted).
Students present evidence of initial iterations (e.g., at least two runs, adjustments, or configurations attempted).

\section{Baseline and Extensions}
\subsection{Baseline Selection and Evaluation (MANDATORY)}
% Rubric: MANDATORY: Students implement and explain their chosen baseline.
MANDATORY: Students implement and explain their chosen baseline.

\subsection{Baseline Reproduction Evidence}
% Rubric: Students show initial results/plots confirming 
% that their baseline results are their own.
Students show initial results/plots confirming that their baseline results are their own.

\subsection{Implemented Extensions/Experiments}
% Rubric: High-level description of extensions/experiments to be conducted.
High-level description of extensions/experiments to be conducted.

\section{Results and Analysis}
% Rubric: Preliminary results or progress metrics encouraged 
% (e.g., sample outputs or early plots).
Preliminary results or progress metrics encouraged (e.g., sample outputs or early plots).

\section{Discussion}
% Rubric: Students briefly discuss challenges so far and potential risks moving forward.
Students briefly discuss challenges so far and potential risks moving forward.

\section{Future Directions}
% Rubric: Planned Next Steps: Students identify 2–3 specific actions to move forward.
Planned Next Steps: Students identify 2–3 specific actions to move forward.

\section{Administrative Details}
\subsection{Team Contributions}
% Rubric: Provide basic role breakdown.
Provide basic role breakdown.

\subsection{Bibliography}
% Rubric: Students include citations for all referenced work (2–3 minimum).
Students include citations for all referenced work (2–3 minimum).

\subsection{GitHub}
% Rubric: Include GitHub link.
Include GitHub link.

\end{document}